\documentclass[10pt,a4paper]{article}
\usepackage[utf8]{inputenc}
\usepackage[spanish]{babel}
\usepackage{url}
\usepackage{hyperref}
\usepackage{amsmath}
\usepackage{amsfonts}
\usepackage{amssymb}
\usepackage{graphicx}
\usepackage{float}
\usepackage{lipsum}
\usepackage{multicol}
\usepackage{xcolor}
\usepackage{tikz}
\include{comandos}
\usepackage[left=2.54cm, right=2.54cm, top=2.54cm, bottom=2.54cm]{geometry}

\title{Mi primer paper}

\begin{document}
	\begin{flushleft}
	    {\small \textbf{Este documento esta hecho en \LaTeX}}
	\end{flushleft}
	
	\vspace{5mm}
	
	\begin{figure}[H]
		\raggedright
		\includegraphics[scale=0.1]{sm.png} \hfill \includegraphics[scale=0.3]{villa.png}
	\end{figure}

	\vspace{7.5mm}
	
	 \tituloDos{Mi segunda plantilla para Overleaf}
	
	\DosAutores{\autor{Manuel Merino$^{1}$}{$^{1}$Universidad Nacillonal Mayor de San Marcos}{Computación Científica}}{\autor{Carla Yauris$^{2}$}{$^{2}$Universidad Nacional Federico Villareal}{Biología}}
	
	 
	
	\dobleLinea
	
	

	\begin{abstract}
		\lipsum[4] \lipsum[5]
		\vspace{2mm}
		\underline{\textbf{Palabras clave:}} \hspace{1mm} \textit{palabra1, palabra2, palabra3.}
	\end{abstract}

	\dobleLinea

	\vspace{5mm}
	
	\begin{multicols}{2}
		\seccion{Introducción}
			\lipsum[2] \lipsum[1] 
			
		\seccion{Conocimientos Previos}
			\lipsum[3] 
	\end{multicols}
	
	\begin{table}[H]
		\centering
		\begin{tabular}{cccc}
			\hline
			\celda{Nombres} & \celda{Apellidos} & \celda{Edad} & \celda{País}\\
			\hline
			\celda{Manuel} & \celda{Merino} & \celda{25} & \celda{Perú}\\
			\hline
			\celda{Pablo} & \celda{Perez} & \celda{27} & \celda{México}\\
			\hline
		\end{tabular}
		\caption{Esta es mi primera tabla en mi paper}
		\label{tb: Tabla1}
	\end{table}

	\begin{multicols}{2}
		\lipsum[1]
		\seccion{Conociendo más de la Inteligencia Artificial.}
			\lipsum
			
			\begin{figure}[H]
				\centering
				\includegraphics[scale=0.1]{sm.png}
				\caption{Esta es mi primera figura del paper}
				\label{fig: Figura1}
			\end{figure}
		
		\lipsum[2]
		
		\seccion{Aplicación y Alcances}
			\lipsum[2]
			
	    \begin{thebibliography}{00}
	    
	    \bibitem{Referencia1}
	        \newblock \textit{Título de la referencia 1}
	        \newblock Autor.
	        \newblock Año y editor.
	    
	    \bibitem{Referencia2}
	        \newblock \textit{Título de la referencia 2}
	        \newblock Autor.
	        \newblock Año y editor.
	        
	    \bibitem{Referencia3}
	        \newblock \textit{Título de la referencia 3}
	        \newblock Autor.
	        \newblock Año y editor.
	        
	    \bibitem{Referencia4}
	        \newblock \textit{Título de la referencia 4}
	        \newblock Autor.
	        \newblock Año y editor.
	    
	    \bibitem{Referencia5}
	        \newblock \textit{Título de la referencia 5}
	        \newblock Autor.
	        \newblock Año y editor.
	        
	    \bibitem{Referencia6}
	        \newblock \textit{Título de la referencia 6}
	        \newblock Autor.
	        \newblock Año y editor.
	    
	    \end{thebibliography}
	\end{multicols}
	
	\vspace{3cm}
	Para un mejor entendimiento de la plantilla ver el video en youtube:\\
	
	\url{https://www.youtube.com/watch?v=34xYuLxnzHo}
	
	
\end{document}